\documentclass[12pt, a4paper]{article}

\usepackage{amsmath, amssymb, amsfonts}
\usepackage{float}
\usepackage{graphicx}
\usepackage{listings}
\usepackage{xcolor}
\usepackage[linesnumbered, ruled]{algorithm2e}
\usepackage{subcaption}
\usepackage[hidelinks]{hyperref}
% Each algorithm is written twice: the algorithm2e block (what the PDF typesets)
% and a plain-text Verbatim block right after it (what the web shows — pandoc
% renders Verbatim as a copyable <pre> code block). \excludecomment makes pdflatex
% SKIP the Verbatim blocks, so the PDF shows only the algorithm2e version; pandoc
% conversely drops each algorithm2e to an (empty, CSS-hidden) div and keeps the
% Verbatim. Because pdflatex never sees the Verbatim, its Unicode maths glyphs
% (⟨⟩ ← τ ∈ ∞ …) need no inputenc/literate setup.
\usepackage{comment}
\excludecomment{Verbatim}

\definecolor{codegreen}{rgb}{0,0.6,0}
\definecolor{codegray}{rgb}{0.5,0.5,0.5}
\definecolor{codepurple}{rgb}{0.58,0,0.82}
\hypersetup{linktoc=all}

\lstset{
    language=Python,
    basicstyle=\ttfamily\small,
    keywordstyle=\color{blue},
    commentstyle=\color{codegreen},
    stringstyle=\color{codepurple},
    breaklines=true,
    numbers=left,
    numberstyle=\tiny\color{codegray},
    frame=single,
    showstringspaces=false
}

\title{AI Masterclass Test 1}
\author{}
\date{}


\begin{document}

\section{Iterative Deepening}
\begin{algorithm}[H]
% \caption{Iterative Deepening (ID)}
\SetKwFunction{Fid}{ID}
\SetKwProg{Fn}{function}{}{end function}

\Fn{\Fid{$\mathcal{P} = \langle S, s_0, A, \tau, c, G \rangle$}}{
    $limit \leftarrow 1$\;
    \While{true}{
        \If{DLS($\mathcal{P}, limit$) == "done"}{
            \tcp*[h]{goal found?} \;
            \Return "done"\;
        }
        \Else{
            $limit \leftarrow limit + 1$\;
        }
    }
}
\end{algorithm}
\begin{Verbatim}
function ID(P = ⟨S, s₀, A, τ, c, G⟩):
    limit ← 1
    while true:
        if DLS(P, limit) == "done":   # goal found?
            return "done"
        else:
            limit ← limit + 1
\end{Verbatim}

\begin{itemize}
    \item If we only perform DLS until depth $d$, but the solution is found at $d+1$, 
    we never find the solution.
    \item Iterative Deepening repeats DLS (calls DLS as subroutine) at increasing depths till solution is found.
    \item We need to regenerate nodes up till depth $d-1$ when we do DLS for depth $d$.
    \item This is tradeoff of time for memory.\newline
    \textbf{Example} $b=10$, $d=5$. BFS would require examining 111,111 nodes, with memory requirement of 100,000 nodes.
    Iterative deepening search 123,456 nodes, with memory requirement of only 50 nodes. Takes 11\% longer with 2000\% reduction in memory.
\end{itemize}

\end{document}
